\ssr{ВВЕДЕНИЕ}

\textbf{Задача коммивояжёра} -- одна из известных задач комбинаторной оптимизации, состоящей в оптимизации замкнутого маршрута по всем заданным точкам. Данная задача имеет длинную историю, впервые ей занялся Эйлер в 1736 году при рассмотрении задачи о семи Кенигсбергских мостах~\cite[с.~641]{eiler}. Первым же неформальным упоминанием задачи считается книга 1832-го года <<Коммивояжер -- как он должен вести себя и что должен делать для того, чтобы доставлять товар и иметь успех в своих делах -- советы старого курьера>>, где была описана проблема, но не дана формальная математическая модель для её решения, на основе которой, в последствие, У.Р. Гамильтон разработал математическую головоломку <<Вокруг света>>~\cite{razvitie}. Первое формальное описание задачи коммивояжёра принадлежит Карлу Менгеру, который математически описал обобщённую задачу коммивояжёра~\cite{history}.

В современном  мире решения задачи коммивояжёра имеют существенное прикладное значение, в частности используясь транспортной логистикой для оптимизации маршрутов доставки товаров~\cite{logistic} или для планирования сложных путешествий через некоторое количество мест~\cite{putesh}, оптимизация конвейера на заводе. Решение задачи коммивояжёра часто используется электронными магазинами, курьерскими службами, например одна из модификаций задачи коммивояжёра легла в основу соревнования E-CUP 2025: ML Challenge~\cite{e-cup}. Количество статей в таких агрегаторах, как \texttt{elibrary.ru}~\cite{elib} и \texttt{google scholar}~\cite{scholar} остаётся на высоком уровне из года в год:
\begin{itemize}
	\item \textbf{2021}: \texttt{elibrary.ru} -- 428, \texttt{google scholar} -- 19300;
	\item \textbf{2022}: \texttt{elibrary.ru} -- 321, \texttt{google scholar} -- 21300;
	\item \textbf{2023}: \texttt{elibrary.ru} -- 478, \texttt{google scholar} -- 21300;
	\item \textbf{2024}: \texttt{elibrary.ru} -- 419, \texttt{google scholar} -- 20700;
	\item \textbf{2025}: \texttt{elibrary.ru} -- 513, \texttt{google scholar} -- 20200.
\end{itemize}
Все это в совокупности доказывает актуальность задачи коммивояжёра.

Несмотря на длительное историю, до сих пор не найден алгоритм, для поиска точного решения задачи со сложностью не выше полиномиальной. В частности полный перебор имеет факториальную сложность. Однако существует множество алгоритмов, позволяющих найти не оптимальное решение, а приближённое, но за обозримое количество операций, например генетический алгоритм, метод имитации отжига, метод градиентного спуска, алгоритм искусственного роя и др. Одним из часто используемых известных эвристических методов является муравьиный алгоритм, однако у него есть ряд существенных проблем, таких как~\cite{impaco}:
\begin{itemize}
	\item медленная сходимость;
	\item высокая вероятность падение в локальный минимум;
	\item подверженность к застою;
	\item сильная зависимость от входных данных и параметров.
\end{itemize}

Эти проблемы частично решаются с помощью модификаций классического алгоритма или комбинацией его с другими известными методами.

Целью данной работы является разработка модификации муравьиного алгоритма с динамическим подбором гетерогенных параметров на основе генетического алгоритма для решения задачи коммивояжёра.
\clearpage
